% !PRINTED 573
\weberpage{573}

%<<CONT
\[
\begin{aligned}
\psi_1(\alpha) &= \psi_3(\alpha^5)\\
\psi_3(\alpha) &= \psi_1(\alpha^3),
\end{aligned}
\]
und wenn man in (8) $\lambda = 1$ setzt
\[
\psi_2(\alpha)\,\psi_3(\alpha) = \psi_1(\alpha)\,\psi_1(\alpha^2);
\]
endlich nach (6)
\[
n = \psi_3(\alpha)\,\psi_3(\alpha^6) = \psi_3(\alpha)\,\psi_1(\alpha^4).
\]
Berechnet man hiernach die Formel \S.\,169, (15)
\[
(\alpha, \eta)^7 = n\,\psi_1(\alpha)\,\psi_2(\alpha)\,\psi_3(\alpha)\,\psi_4(\alpha)\,\psi_5(\alpha),
\]
so ergiebt sich
\weq{9}{(\alpha, \eta)^7 = \psi_1(\alpha)^4\,\psi_1(\alpha^2)^2\,\psi_1(\alpha^4).}
Durch diese Formel l\"asst sich unmittelbar die Relation aus \S.\,169, (14)
\[
(\alpha, \eta)^{14} = (\alpha^2, \eta)^7\,\psi_1(\alpha)^7
\]
verificiren.

Wir wollen diese Betrachtungen \"uber die Functionen $\psi$ mit dem Beweise eines Satzes beschliessen, der uns sp\"ater n\"utzlich sein wird.

Bestimmen wir eine positive ganze Zahl $\nu < m$ aus der Congruenz
\weq{10}{\lambda + \mu + \nu \equiv 0 \pmod{m},}
so k\"onnen wir die Function $\psi_{\lambda,\mu}(\varepsilon)$, wie sie durch \S.\,169, (6) definirt ist, so darstellen
\weq{11}{\psi_{\lambda,\mu}(\varepsilon) = \sum_{1,\,n-2}^{t} \varepsilon^{\mu\,\mathrm{ind}\,t + \nu\,\mathrm{ind}(t+1)}.}
Hierin wollen wir nun die $m^{\text{te}}$ Einheitswurzel $\varepsilon$ durch die primitive Congruenzwurzel $g$ der Primzahl $n$, die dem Indexsystem zu Grunde liegt, ersetzen, d.\,h. wir wollen die ganze Zahl
\weq{12}{\psi_{\lambda,\mu}(g) = \sum_{1,\,n-2}^{t} g^{\mu\,\mathrm{ind}\,t + \nu\,\mathrm{ind}(t+1)} \equiv \sum_{1,\,n-2}^{t} t^\mu(t+1)^\nu \pmod{n}}
bilden. Betrachten wir diese Zahl nach dem Modul $n$, so k\"onnen wir die Summation bis $t = n - 1$ ausdehnen, da f\"ur diesen Werth $(t + 1)$ congruent mit Null ist. Wir erhalten also
\weq{13}{\psi_{\lambda,\mu}(g) \equiv \sum_{1,\,n-1}^{t} t^\mu(t+1)^\nu \pmod{n}.}
Auf die Potenz $(t + 1)^\nu$ k\"onnen wir den binomischen Lehrsatz anwenden und erhalten, wenn, wie fr\"uher, $B_h^{(\nu)}$ die Binomialco\"efficienten bedeuten,
\weq{14}{\psi_{\lambda,\mu}(g) \equiv \sum_{0,\,\nu}^{h} B_h^{(\nu)} \sum_{1,\,n-1}^{t} t^{\mu+h} \pmod{n}.}
